\section{Motivación}

\begin{frame}{Contexto}

  \begin{itemize}
    \item El Aprendizaje por Refuerzo Profundo ha mostrado resultados sobresalientes y prometedores en areas como:
          \begin{itemize}
            \item Tareas de control: robotica
            \item Juegos Complejos: Atari
            \item Entrenamiento de LLMs: fine tunning de chatGPT.
          \end{itemize}
    \item Predomina el uso de arquitecturas neuronales tradicionales:
          \begin{itemize}
            \item MLP
            \item CNN
            \item RNN
          \end{itemize}
  \end{itemize}

  \begin{block}{}
    ¿Que ocurre si utilizamos arquitecturas inspiradas biologicamente?
  \end{block}
\end{frame}

\begin{frame}{Antecedentes}  
  Trabajos anteriores buscaron responder dicha pregunta y lograron:
  \begin{itemize}
      \item Integrar una \textbf{arquitectura neuronal bioinspirada} (Tap Withdrawal Circuit de \textit{C. elegans}) en el marco del aprendizaje por refuerzo.
      \item Plantear nuevos modelos de dinamica neuronal (\textbf{modelo Fiuri}) y compararlos con modelos estándar (IandF, Izikievich).
      \item Utilizar metodos de entrenamiento alternativos al descenso por el gradiente (Optimizacion bayesiana, algoritmos geneticos y busquedas aleatoreas).
  \end{itemize}
  
  Nos preguntamos si es posible:
  \begin{itemize}
    \item Implementar de manera eficiente y vectorizada
    \item Permitir el uso del el descenso por el gradiente
    \item Integrar a los algoritmos de Deep RL estandar
  \end{itemize}
\end{frame}

\begin{frame}{Objetivo y Desafios}
  \begin{block}{Objetivo}
    Integrar la arquitectura \textbf{TWC + FIU} al framework de PyTorch e intentar resolver un problema de control continuo utilizando un algoritmo estandar Deep RL.
  \end{block}

  \textbf{Desafios}
  \begin{itemize}
    \item Las arquitecturas bioinspiradas presentan:
    \begin{itemize}
        \item Dinámica interna recurrente
        \item No linealidades inspiradas biológicamente
        \item Posibles problemas de diferenciabilidad
    \end{itemize}
    
    \item La mayoría de los trabajos previos:
    \begin{itemize}
        \item Evitan el descenso por gradiente
        \item Utilizan métodos evolutivos o búsqueda directa en el espacio de políticas
    \end{itemize}
\end{itemize}
\end{frame}